% ============================================================
% Chapter 9 — The rotating ring
%
% Original work of this thesis, reported here for the first time; there is
% no external source to cite for the derivations, the diagnostics, or the
% measurements.
% Numerical values are those of research/experiment1-rotating-ring/data.
% ============================================================

\chapter{The rotating ring: a two-dimensional dynamic object}
\label{ch:ring}

Chapters~\ref{ch:gaussian} and~\ref{ch:laplace} studied a chain of scalar
frames, where the joint score is either exactly Gaussian or a controlled
departure from it. Chapter~\ref{ch:development} set aside a second family of
dynamic objects --- planar trajectories near a circle --- because they are
nonlinear in the coordinates in which the diffusion acts. This chapter takes
up the first of them. It shows that a
Cartesian surrogate of the model is exactly solvable, that the model as drawn
is solvable numerically to a measured accuracy, and that both admit the same
structural description of the joint score.

This chapter is a self-contained secondary contribution. It does not feed the
chain results of Chapters~\ref{ch:gaussian}--\ref{ch:em-results} and nothing
there depends on it; it is included because it settles, exactly and in a model
where the calculation can be done, a question the chain model raises but cannot
answer --- whether per-frame marginals can carry the information that
distinguishes one dynamics from another. Chapter~\ref{ch:development} settles
this for the stationary Gaussian chain, where the marginals are
$\alpha$-free outright (Proposition~\ref{prop:dev-marginal}), and shows that the
same argument does not extend to non-Gaussian innovations. Here it becomes a
theorem of a different kind: for one specific dynamical parameter, in a model
whose symmetry makes the calculation exact, the one-frame marginals carry
\emph{zero} Fisher information while the trajectory carries a positive amount.

The question the chapter answers is narrower than the models it uses.

\begin{quote}
How can the Markov structure of the clean rotating trajectory be exploited to
compute and interpret the joint diffusion score?
\end{quote}

Four measurable subquestions make this operational: how far the joint score
differs from a collection of one-frame marginal scores; how strongly the
observation at frame $j$ influences the score block at frame $k$; how that
influence decays with the temporal lag $|j-k|$; and how many neighbouring
frames are needed to reproduce the score block to a stated accuracy. Every
answer below is a statement about one model at one parameter set, in the sense
fixed in Section~\ref{sec:rq}.

\section{Model and notation}\label{sec:ring-model}

\subsection{The clean process}

The clean state is a radius and a phase, evolving in the internal time $u$ as
two independent diffusions,
\begin{align}
  \dd R_u &= -\kappa (R_u - 1)\,\dd u + \sqrt{2 D_r}\, \dd B_u^{(r)} ,
  \label{eq:ring-radial-sde} \\
  \dd \Theta_u &= \omega\, \dd u + \sqrt{2 D_\theta}\, \dd B_u^{(\theta)} ,
  \label{eq:ring-angular-sde}
\end{align}
observed through the polar embedding
\begin{equation}
  a_u = R_u \big(\cos\Theta_u,\ \sin\Theta_u\big) \in \R^2 .
  \label{eq:ring-embedding}
\end{equation}
We use the simplest constant-angular-drift reading of the model;
Section~\ref{sec:ring-limits} records the alternatives.

A point of $\R^2$ is described by its distance from the origin and the angle it
makes with the first axis, and at each point the plane carries the orthonormal
moving frame
\begin{equation}
  e_r = (\cos\theta,\ \sin\theta), \qquad
  e_\theta = (-\sin\theta,\ \cos\theta) ,
  \label{eq:ring-moving-frame}
\end{equation}
in which $e_\theta$ is $e_r$ turned by a quarter revolution. Moving along $e_r$
changes the radius at fixed angle, moving along $e_\theta$ changes the angle at
fixed radius, and the area element is $\dd x\,\dd y = r\,\dd r\,\dd\theta$. The
decomposition matters because the two directions are treated very differently
by \eqref{eq:ring-radial-sde} and \eqref{eq:ring-angular-sde}: the radius is
pulled back towards one, the phase drifts and diffuses freely. Measuring the
consequence of that asymmetry for the score is one purpose of the chapter.

\subsection{The radial drift is a confining force}

Equation~\eqref{eq:ring-radial-sde} is the overdamped Langevin equation
\eqref{eq:langevin-sde} for the quadratic potential
\begin{equation}
  V(r) = \tfrac{\kappa}{2}\,(r-1)^2 .
  \label{eq:ring-potential}
\end{equation}
A potential stores energy, and the force it generates is minus its slope,
\begin{equation}
  F(r) = -V'(r) = -\kappa\,(r-1) .
  \label{eq:ring-force}
\end{equation}
The sign carries the content: a system driven by $-V'$ moves downhill in $V$,
since $\tfrac{\dd}{\dd u} V(r_u) = V'(r)\,\dot r = -V'(r)^2 \le 0$. Reading the
two cases of \eqref{eq:ring-force}, if $r > 1$ then $F < 0$ and the force points
inward, and if $r < 1$ then $F > 0$ and it points outward. The radius is driven
towards one from either side at a rate proportional to the displacement, with
$\kappa$ setting the stiffness. By \eqref{eq:langevin-boltzmann} the stationary
radial law is $\Nd(1, D_r/\kappa)$.

\subsection{One trajectory is one data point}

The trajectory is sampled at internal times $u_k = k h$ for
$k = 0, \dots, K-1$, and the complete datum is the stacked vector
\begin{equation}
  a = (a_0, a_1, \dots, a_{K-1}) \in \R^{2K} .
  \label{eq:ring-datum}
\end{equation}
The $K$ frames are correlated coordinates of one data point, not $K$ samples.
Relative to Chapter~\ref{ch:gaussian} the only change of setting is that each
frame is now two-dimensional, so the score has $K$ blocks of size two rather
than $K$ scalar components; the forward channel is unchanged,
\begin{equation}
  x = e^{-t} a + \sqrt{\Deltat}\, \xi , \qquad \xi \sim \Nd(0, I_{2K}) ,
  \label{eq:ring-channel}
\end{equation}
which is the channel of Section~\ref{sec:ou} acting on all $2K$ coordinates at
once. Figure~\ref{fig:ring-model} shows the model and the two clocks: internal
time orders the frames along the path, diffusion time corrupts them
simultaneously.

\begin{figure}[t]\centering
\includegraphics[width=\linewidth]{fig_ring_model.pdf}
\caption{The rotating-ring model, at the parameters used throughout this
chapter ($K = 20$, $\kappa = 2$, $D_r = 0.015$, $D_\theta = 0.005$,
$\omega = 1$, $h = 2\pi/20$). (a) Polar coordinates and the moving frame
\eqref{eq:ring-moving-frame}. (b) The potential \eqref{eq:ring-potential} and
the force \eqref{eq:ring-force} it generates; the force changes sign at
$r = 1$, which is therefore the stable radius. (c) One clean trajectory, a
noisy coherently turning path near the unit circle; the arrow marks the sense
of rotation. (d) The same trajectory after the channel
\eqref{eq:ring-channel} at $t = 0.5$, from a single standardised noise
realisation.}
\label{fig:ring-model}
\end{figure}

\section{The gauge, and the exact score}\label{sec:ring-surrogate}

\subsection{The surrogate, and what it keeps}

The surrogate retains a ring, a coherent rotation, and dependence along the
trajectory, and drops the rest:
\begin{align}
  z_0 &\sim p_\lambda(z) \propto
        \exp\!\Big[-\frac{(\norm{z}-1)^2}{2\lambda}\Big] ,
  \label{eq:ring-sur-prior} \\
  z_{k+1} &= R_\psi z_k + \sigma \eta_{k+1} ,
        \qquad \eta_{k+1} \stackrel{\text{iid}}{\sim} \Nd(0, I_2) ,
  \label{eq:ring-sur-dynamics}
\end{align}
with $R_\psi$ the planar rotation through the one-frame angle
$\psi = \omega h$.

\begin{remark}[The surrogate is not a discretisation of the polar model]
\label{rem:ring-surrogate-scope}
In \eqref{eq:ring-sur-dynamics} the ring acts only on $z_0$. Later frames form
a rotating Cartesian random walk and their radial variance grows without bound,
whereas in \eqref{eq:ring-radial-sde} a confining force acts at every internal
time. The surrogate is a solvable companion whose role is to exhibit the
mechanism in closed form, not to approximate the model of
Section~\ref{sec:ring-model}.
\end{remark}

\subsection{The known rotation is an orthogonal change of frame}

The rotation is the object to be removed, so we write it out:
\begin{equation}
  R_\psi =
  \begin{pmatrix} \cos\psi & -\sin\psi \\ \sin\psi & \cos\psi \end{pmatrix} .
  \label{eq:ring-rotation}
\end{equation}
Its columns are the standard basis vectors, each turned through $\psi$. They
are unit vectors, since $\cos^2\psi + \sin^2\psi = 1$, and orthogonal, since
$\cos\psi(-\sin\psi) + \sin\psi\cos\psi = 0$. In matrix form,
\begin{equation}
  R_\psi^{\T} R_\psi = I_2, \qquad \det R_\psi = 1, \qquad
  R_\psi^{\T} = R_{-\psi}, \qquad R_\psi R_\varphi = R_{\psi+\varphi} ,
  \label{eq:ring-rot-props}
\end{equation}
and orthogonality gives
$\norm{R_\psi z}^2 = z^{\T} R_\psi^{\T} R_\psi z = \norm{z}^2$: lengths and
inner products are preserved.

\begin{lemma}[Rotations preserve the standard isotropic Gaussian]
\label{lem:ring-isotropy}
If $\eta \sim \Nd(0, I_2)$ and $R$ is orthogonal, then $R\eta \sim \Nd(0,I_2)$.
\end{lemma}

\begin{proof}
Appendix~\ref{app:ring}.
\end{proof}

The isotropy is essential. A rotation would not preserve $\Nd(0,\Sigma)$ for
general $\Sigma$, and the reduction below is exact only because both the
innovation in \eqref{eq:ring-sur-dynamics} and the channel
\eqref{eq:ring-channel} are isotropic.

\begin{lemma}[Gauge reduction]\label{lem:ring-gauge}
Define the de-rotated frame $y_k = R_{-k\psi} z_k$ and stack the per-frame
rotations into
$U_\psi = \mathrm{diag}(I_2, R_{-\psi}, \dots, R_{-(K-1)\psi})$. Then
\begin{equation}
  y_{k+1} = y_k + \sigma \widetilde\eta_{k+1} ,
  \qquad \widetilde\eta_{k+1} \stackrel{\text{iid}}{\sim} \Nd(0, I_2) ,
  \label{eq:ring-gauged-rw}
\end{equation}
$U_\psi$ is orthogonal, and the noised density and score satisfy
\begin{equation}
  P_t(z \,|\, \psi) = P^{(0)}_t(U_\psi z) ,
  \qquad
  \score_\psi(z,t) = U_\psi^{\T}\, \score^{(0)}(U_\psi z, t) ,
  \label{eq:ring-gauge-score}
\end{equation}
where the superscript zero denotes the quantities of the non-rotating model.
\end{lemma}

\begin{proof}
Appendix~\ref{app:ring}.
\end{proof}

The term \emph{gauge} is used in the sense standard in physics, of a change of
description that leaves the content untouched and may be chosen to simplify the
equations. Equation~\eqref{eq:ring-gauge-score} says that a known rotation adds
no statistical difficulty: de-rotate the data, solve the non-rotating problem,
rotate the score back. Since $U_\psi^{\T} = U_\psi^{-1}$, the operation is
exactly invertible and destroys no information. It would not hold for an
\emph{unknown} $\psi$, which is a genuine latent variable
(Section~\ref{sec:ring-limits}). Figure~\ref{fig:ring-gauge} shows the two
frames.

\begin{figure}[t]\centering
\includegraphics[width=\linewidth]{fig_ring_gauge.pdf}
\caption{The rotation and the gauge. (a) $R_\psi$ turns every vector through
the same angle without changing its length, which is the statement
$R_\psi^{\T}R_\psi = I_2$. (b) Three surrogate trajectories in the laboratory
frame ($K = 14$, $\sigma = 0.16$, $\psi = 2\pi/14$; the marked point is
$k = 0$), coherently circulating. (c) The same trajectories after $U_\psi$: a
random walk anchored on the ring. Because $U_\psi$ is orthogonal the reduction
is exact at every diffusion time; the numerical round trip agrees to machine
precision (Table~\ref{tab:ring-validation}).}
\label{fig:ring-gauge}
\end{figure}

\subsection{The clean law after the gauge}

Iterating \eqref{eq:ring-gauged-rw} and cancelling the intermediate terms gives
the anchored representation \eqref{eq:ring-unrolled}, whose derivation, together
with the second-difference identity it satisfies, is in
Appendix~\ref{app:ring}. The identity is verified to machine zero
(Table~\ref{tab:ring-validation}).
\subsection{The noised law and the exact score}

Conditioning on the anchor, \eqref{eq:ring-unrolled} gives
$\E[y_i \,|\, A = \alpha] = \alpha$ and
\begin{equation}
  \mathrm{Cov}(y_i, y_j \,|\, A)
  = \sigma^2\, \E\Big[\Big(\textstyle\sum_{r \le i}\widetilde\eta_r\Big)
      \Big(\textstyle\sum_{s \le j}\widetilde\eta_s\Big)^{\!\T}\Big]
  = \sigma^2 \min(i,j)\, I_2 ,
  \label{eq:ring-K}
\end{equation}
since $\E[\widetilde\eta_r \widetilde\eta_s^{\T}] = \delta_{rs} I_2$ leaves only
the $\min(i,j)$ indices common to both sums. Compactly,
$y \,|\, A = \alpha \sim \Nd(\mathbf{1} \otimes \alpha,\ G \otimes I_2)$
with $G_{ij} = \sigma^2\min(i,j)$, where $\mathbf 1 \in \R^K$ is the
all-ones vector. The two Kronecker products separate the temporal index from
the spatial one: $\mathbf 1 \otimes \alpha$ repeats the same two-dimensional
anchor in every frame, and
$(G \otimes I_2)_{(i,\alpha'),(j,\beta')}
= G_{ij}\,\delta_{\alpha'\beta'}$, so $G$ carries all the temporal
correlation while $I_2$ states that the two spatial coordinates are
uncorrelated and share the same temporal covariance. Applying the channel and
using $I_{2K} = I_K \otimes I_2$ to combine the covariances,
\begin{equation}
  x \,|\, A = \alpha
  \ \sim\ \Nd\big(e^{-t}\,\mathbf 1 \otimes \alpha,\ C_t \otimes I_2\big),
  \qquad
  C_t = e^{-2t}G + \Deltat\, I_K ,
  \label{eq:ring-Ct}
\end{equation}
in which $C_t$ does not depend on $\alpha$: the noised law is a two-parameter
location mixture of one fixed Gaussian.

\begin{proposition}[Exact surrogate score and response]\label{prop:ring-exact}
Let $q = C_t^{-1}\mathbf 1$, $\beta = e^{-2t}\mathbf 1^{\T} C_t^{-1}\mathbf 1$
and $h = e^{-t}\sum_j q_j x_j \in \R^2$. Then the anchor posterior is
\begin{equation}
  p(\alpha \,|\, x) \propto
  \exp\!\Big[ -\frac{(\norm{\alpha}-1)^2}{2\lambda}
              - \frac{\beta}{2}\norm{\alpha}^2 + h^{\T}\alpha \Big] ,
  \label{eq:ring-anchor-posterior}
\end{equation}
and the score and response blocks are
\begin{equation}
  \score_k(x,t) = -\big(C_t^{-1}x\big)_k + e^{-t} q_k\, \E[A \,|\, x] ,
  \qquad
  J_{kj}(x,t) = -\big(C_t^{-1}\big)_{kj} I_2
                + e^{-2t}\, q_k q_j\, \mathrm{Cov}(A \,|\, x) .
  \label{eq:ring-surrogate-score}
\end{equation}
\end{proposition}

\begin{proof}
Appendix~\ref{app:ring}.
\end{proof}

The $2K$-dimensional integral over clean trajectories has collapsed to a
two-dimensional posterior over the shared anchor, which is the whole
simplification and a direct consequence of \eqref{eq:ring-unrolled}. Its two
moments are cheap: the prior and the $\norm{\alpha}^2$ term are rotationally
symmetric, so the posterior depends on the direction of $\alpha$ only through
$h^{\T}\alpha$, and
$\E[A\,|\,x] = \mu(\norm{h})\,\widehat h$ with
$\mathrm{Cov}(A\,|\,x) = v_\perp I_2 + (v_\parallel - v_\perp)
\widehat h \widehat h^{\T}$, three scalar functions of one scalar argument,
tabulated once per diffusion time by one-dimensional quadrature.

The structural content of \eqref{eq:ring-surrogate-score} is that the temporal
coefficient of the second term is the outer product $q q^{\T}$, which is of
rank one. The nonlocal part of the response is therefore not a chain of local
couplings but a single shared latent variable seen by every frame.
Figure~\ref{fig:ring-surrogate-response} separates the two terms.

\begin{figure}[t]\centering
\includegraphics[width=\linewidth]{fig_ring_surrogate_response.pdf}
\caption{Exact decomposition of the surrogate response
\eqref{eq:ring-surrogate-score} at $t = 0.20$, $K = 30$, $\lambda = 0.05$,
$\sigma = 0.15$, averaged over $250$ noisy trajectories, each panel normalised
by its own maximum. (a) The chain term $|C_t^{-1}|$ is concentrated near the
diagonal. (b) The anchor coefficient $|qq^{\T}|$ is rank one and spread over
all lags. (c) The mean norm of the complete $2\times2$ blocks. The analytic
Jacobian agrees with central finite differences to $9.2\times10^{-6}$
(Table~\ref{tab:ring-validation}).}
\label{fig:ring-surrogate-response}
\end{figure}

\begin{proposition}[What the surrogate establishes]
\label{prop:ring-surrogate-conclusion}
Clean Markov locality does not survive marginalisation through the OU channel
as locality of the noised score: by \eqref{eq:ring-surrogate-score} the
response has nonzero blocks at every lag for every $t > 0$. The nonlocal part
is nonetheless structured, decomposing exactly into a near-diagonal chain term
and a rank-one correction generated by the shared ring anchor.
\end{proposition}

\section{Numerical validation}\label{sec:ring-validation}

Every identity of this chapter is checked against an independently computed
reference, in the protocol of Appendix~\ref{app:repro}.

\begin{table}[t]\centering
\caption{Validation of the rotating-ring computations. The response identity
and the analytic surrogate Jacobian are compared with central finite
differences of the score; the grid check compares two polar resolutions.}
\label{tab:ring-validation}
\begin{tabular}{@{}lrr@{}}
\toprule
check & discrepancy & tolerance \\
\midrule
posterior response identity vs finite differences
  & $6.63\times10^{-9}$ & $10^{-7}$ \\
analytic surrogate Jacobian \eqref{eq:ring-surrogate-score} vs finite differences
  & $9.18\times10^{-6}$ & $10^{-4}$ \\
$41\times128$ against $51\times192$ polar response grid
  & $3.04\times10^{-11}$ & $10^{-7}$ \\
rotation gauge round trip (Lemma~\ref{lem:ring-gauge})
  & $2.22\times10^{-16}$ & $10^{-12}$ \\
clean score on an exact co-rotating ring path
  & $0$ (exact) & $10^{-12}$ \\
\bottomrule
\end{tabular}
\end{table}


Beyond Table~\ref{tab:ring-validation} the checks include row normalisation of
the transition kernels, global rotation equivariance of the score, stability
under halving the finite-difference step, and the exact vanishing of the clean
surrogate score on a perfect co-rotating ring path. One caveat on grid
convergence: the difference between the $41\times128$ and $51\times192$ grids
is small because the baseline likelihoods and kernels are already smooth at
those resolutions, and a deliberately coarse $11\times24$ grid does show
visible deviations. The production grid is converged for these parameters
rather than trivially grid-independent.

\section{Marginal blindness}\label{sec:ring-blindness}

Everything above concerns how far influence \emph{reaches} in the noised score.
This section establishes something stronger and of a different kind: there is a
parameter of this model about which the per-frame marginals carry not less
information, but none.

\begin{theorem}[Marginal blindness]\label{thm:ring-blindness}
For the rotating ring with uniform initial phase, every single-frame marginal
is independent of the rotation $\psi$, at every noise level. Consequently a
model built from per-frame marginals carries exactly zero Fisher information
about the rotation.
\end{theorem}

\begin{proof}
Unrolling the recursion,
$z_u = R_\psi^{\,u} z_0 + \sigma\sum_{k\le u} R_\psi^{\,u-k}\eta_k$. Each
$\eta_k$ is isotropic Gaussian, and a rotation of an isotropic Gaussian has the
same law, so $R_\psi^{\,u-k}\eta_k \stackrel{d}{=} \eta_k$ independently across $k$. The
initial phase $\theta_0$ is uniform on the circle, so $z_0$ has a rotationally
invariant law and $R_\psi^{\,u} z_0 \stackrel{d}{=} z_0$. The two terms are independent,
hence
\begin{equation}
  z_u \;\stackrel{d}{=}\; z_0 + \sigma\sqrt{u}\,\xi,
  \qquad \xi \sim \mathcal{N}(0, I_2),
  \label{eq:ring-blind-law}
\end{equation}
in which $\psi$ does not appear. The channel acts coordinatewise and
independently of $\psi$, so the noised frame is likewise $\psi$-free. Therefore
$\partial_\psi \log P_t(x_u) \equiv 0$ for every frame and every noise level,
and a sum of per-frame log-densities cannot contain information no term
contains.
\end{proof}

\begin{remark}[What the theorem does not say]
The marginals are not uninformative in general. They determine the width and
centre of the confining well and the walk noise perfectly well, and by
\eqref{eq:ring-blind-law} the per-frame law is a visible, non-Gaussian ring
whose radius grows as $\sqrt{u}$. They are blind to the \emph{rotation}
specifically. Nor is the joint easy: the joint carries the information, and
extracting it is what the estimator does.

The sharpness is what matters. A per-frame model is not less informative about
the dynamics, nor informative only at low noise. It carries zero, exactly, at
every noise level. Radii are not evidence.
\end{remark}

\section{Three estimation targets}\label{sec:ring-estimation}

The theorem furnishes an unusually clean experiment. This model carries a
parameter the marginals can see --- the confining potential
$V(r) = (r-r_*)^2/2\lambda$ --- and a parameter they provably cannot, so
running both through the same estimator isolates the effect of the structure
rather than of the method. Three targets, each in two arms differing only in
which likelihood scores a trajectory:

\begin{center}\small
\begin{tabular}{@{}llll@{}}
\toprule
 & target & free parameters & visible to marginals? \\
\midrule
(a) & the potential $(\lambda, r_*)$ & two & \textbf{yes} \\
(b) & a single rotation $\psi$ & one & \textbf{no} \\
(c) & a law over rotations $p(\psi)$ & a distribution & \textbf{no} \\
\bottomrule
\end{tabular}
\end{center}

\noindent Target~(a) defines the contrast rather than supplying a result: it is
the marginal-visible control against which (b) and (c) are read, and it is what
makes the design isolate the structure rather than the estimator. \textbf{No
measurement of (a) is reported in this thesis.} What follows is (b) and (c).

For (c) the estimator is expectation--maximisation over a grid of angles, with
responsibilities $r_\mu(\psi)\propto p(\psi)P_t(x^\mu\mid\psi)$ and the
update $p(\psi)\leftarrow M^{-1}\sum_\mu r_\mu(\psi)$. The same sufficiency
that makes the expected-transition statistic a one-shot quantity on the chain
(Chapter~\ref{ch:em-method}) appears here: the $\psi$-conditional
log-likelihood does not depend on $p(\psi)$, so the matrix of log-densities is
computed once and every iteration reweights it.

The blind arm is not a baseline chosen to lose. It is the theorem's own
control: its likelihood is constant in $\psi$, so its responsibilities equal
$p(\psi)$, its update is the identity, and $p(\psi)$ cannot leave its uniform
initialisation. If it moves, the implementation is wrong rather than the
theorem.

\paragraph{The blind arm returns exactly the null.} Over sixteen seeds, seven
budgets and twelve noise levels --- $1344$ cells --- the per-frame arm's
estimate moves $0.0000$ in total variation from its uniform initialisation,
and over the $480$ cells of target (b) its profile likelihood in $\psi$ has
range $0.00\times10^{0}$ nats. This is a machine-precision zero rather than a
small number, which is what a theorem predicting exactly no information
requires.

It is worth being precise about what that zero does and does not evidence. The
per-frame likelihood is implemented by a function that takes no $\psi$
argument, so its independence of $\psi$ is a property of the code path, not a
measurement: the zero confirms that the estimator is wired to the marginal it
is supposed to use, and would catch a leak, but it cannot corroborate the
theorem against an independently constructed alternative.

The independent check is a paired one. Trajectories are simulated separately at
$\psi \in \{0,\,\pi/6,\,\pi/3,\,1.1\}$, sharing no code with the estimator, and
two statistics are computed on the \emph{samples}. As a negative control, the
energy distance between one-frame clouds at different $\psi$ is compared
against the same distance between two independent draws at \emph{equal} $\psi$,
which is the sampling floor: the floor is $1.5\times10^{-2}$ and the
cross-$\psi$ distances are $2.6$ to $4.7\times10^{-3}$ --- strictly
\emph{below} the noise of having drawn twice, both on clean trajectories and
after OU noising at $t=0.5$. As a positive control, the lag-one cross-covariance
of those same samples must equal $s^2 R_\psi$, so its antisymmetric part is
$\sin\psi$; measured, it is $-0.001,\,0.503,\,0.865,\,0.893$ against
$\sin\psi = 0,\,0.500,\,0.866,\,0.891$.

The pairing is the point. A null result on its own cannot distinguish a true
invariance from a statistic too blunt to detect anything; here the same
samples, at the same size, move by $0.9$ in the joint statistic and by less
than their own resampling noise in the marginal one. Both controls are in
\codefile{tests/test\_ring\_blindness\_controls.py}.

\paragraph{The joint arm recovers the rotation, and loses it to the channel.}
Estimating a single angle, the error is $2.1\pm0.7^\circ$ pooled over the whole
schedule. Recovering the two-species law
$p^*(\psi)=\tfrac12\delta_{+\omega}+\tfrac12\delta_{-\omega}$ against a truth
of $\omega = 30^\circ$ and a null of $90^\circ$:

\begin{center}\small
\begin{tabular}{@{}lcccccc@{}}
\toprule
$t$ & $0.05$ & $0.32$ & $0.68$ & $0.98$ & $1.43$ & $3.00$ \\
\midrule
median $|\hat\omega - \omega|$ & $0.15^\circ$ & $0.24^\circ$ & $0.90^\circ$
& $2.86^\circ$ & $7.09^\circ$ & $58.8^\circ$ \\
\bottomrule
\end{tabular}
\end{center}

Information about the rotation survives until $t\approx1$ and is gone by
$t=3$, where the estimate has returned to the null. That is the channel doing
what a channel does, and it is legible here in a way it is not on the chain,
because the null is known exactly rather than estimated.

\begin{remark}[A normalisation that fails silently]
The radial density of the confining well is naturally written unnormalised,
and its normaliser depends on both $\lambda$ and $r_*$. While only $\psi$ is
estimated the term is constant and cancels in every responsibility --- which
is why it can be absent for a long time without anything appearing wrong. The
moment $\lambda$ is estimated it stops being constant: a wider well contains
more mass, the likelihood grows without bound in $\lambda$, and the estimate
pins to the top of whatever grid it is given. The failure produces a smooth,
confident, wrong answer rather than an error. It is recorded here because the
same shape of mistake --- a quantity that is constant in the case one has been
testing, and is not constant in the case one moves to --- is the most common
way a correct derivation becomes an incorrect measurement.
\end{remark}

\section{What this chapter establishes}\label{sec:ring-answers}

Three things, in increasing order of how much they constrain what a method can
do.

First, the rotation is a \emph{gauge}. Because $U_\psi$ is orthogonal it
commutes with the isotropic channel, so at known $\psi$ the model reduces to
$\psi = 0$ and the conditional density is closed form at any trajectory length
--- a quadratic form plus a single one-dimensional radial integral. There is a
ground truth here, computed rather than estimated, and it is what everything
else is measured against.

Second, the clean law is local but the noised score is not. For $t>0$,
$\score_k(x,t)$ is not a function of $x_{k-1},x_k,x_{k+1}$ alone: integrating
out uncertain latent states creates posterior correlations at every lag, and
the anchor term of \eqref{eq:ring-surrogate-score} is rank one but touches
every pair of frames. Clean Markovity buys a local \emph{generative} law, not a
local score.

Third, and this is the result that does not have an analogue on the chain,
there is a parameter of the model about which per-frame marginals carry
\emph{exactly zero} information (Theorem~\ref{thm:ring-blindness}), and it is
nonetheless recoverable from the joint. The blind arm of
\S\ref{sec:ring-estimation} does not merely do worse; it cannot move off its
initialisation, and over $1{,}824$ cells it does not, to machine precision.
That is a sharper statement than a comparison of error rates, because the null
is known exactly rather than estimated.

\begin{table}[t]\centering
\caption{Status of the results of this chapter.}
\label{tab:ring-claims}
\begin{tabular}{@{}p{0.50\linewidth}p{0.46\linewidth}@{}}
\toprule
result & status \\
\midrule
Gauge reduction and score transformation
  (Lemma~\ref{lem:ring-gauge}) & exact \\
Exact surrogate score and rank-one response
  (Proposition~\ref{prop:ring-exact}) & exact \\
Marginal blindness (Theorem~\ref{thm:ring-blindness}) & exact \\
Recovery of $p(\psi)$ from the joint, and its decay with $t$
  & numerical, sixteen seeds \\
The blind arm returning the null to machine precision
  & numerical, $1{,}824$ cells; a consistency check on the estimator's
    wiring, not on the theorem \\
Paired sample-level controls: one-frame law invariant to $\psi$ below the
  resampling floor, lag-one joint tracking $\sin\psi$
  & numerical, independent of the estimator \\
Recovery of the confining potential $(\lambda, r_*)$
  & not reported here (defines the control, see
    \S\ref{sec:ring-estimation}) \\
Behaviour in $K$, and inference of the phase rather than its use
  & open \\
\bottomrule
\end{tabular}
\end{table}

\section{Limitations}\label{sec:ring-limits}

\begin{enumerate}
  \item \textbf{One parameter set for the geometry.} The ring width, walk noise
  and trajectory length are held at a single baseline. Only the noise level and
  the data budget are swept, and nothing is claimed about scaling in $K$.
  \item \textbf{Radius positivity.} The radial process is formally supported on
  all of $\R$. At the baseline the probability of a negative radius is
  negligible, but a strict model would reflect at the origin or use a positive
  diffusion.
  \item \textbf{Known rotation in the gauge argument.}
  Lemma~\ref{lem:ring-gauge} assumes $\psi$ known; \S\ref{sec:ring-estimation}
  is what happens when it is not, and the two-species case is the smallest
  version of the speciation phenomenon discussed in
  Section~\ref{sec:statphys-diffusion}.
  \item \textbf{Grid approximation.} The radial integral is numerically
  converged at these parameters but is a quadrature, not a symbolic closed
  form. The validation of \S\ref{sec:ring-validation} bounds the error; it does
  not remove it.
  \item \textbf{The marginal-visible control is not measured.} Target~(a) of
  \S\ref{sec:ring-estimation} specifies the contrast but no result for it is
  reported, so the claim that the two targets behave differently rests on the
  theorem rather than on a matched pair of measurements. Nothing else in the
  chapter depends on it.
\end{enumerate}

The immediate extensions are to vary $K$ systematically, to infer the angular

phase rather than use an oracle basis, to add an unknown rotation class, and to
test whether a time-dependent local architecture with receptive field
$L_\epsilon(t)$ can match the full score.
